% Theory section for: "Does random maneuvering help a drone slip past sensors?
% A diffusion model for counter-UAS sensor layouts."
% Standalone draft: compiles on its own; paste the body into the manuscript later.
%
% PROVENANCE (remove before submission, but keep the disclosure the journal requires):
%  - Most proofs and the asymptotic expansions were derived with help from AI assistants;
%    the arguments were checked step by step and every number was reproduced in code
%    (move the check scripts into experiments/ before submission).               [TODO]
%  - Prop 9 (local optimality) is a computer-assisted result produced externally; it must be
%    re-run from its published code before it is cited as proved.        [TODO]
%  - Prop 10 (branch and bound) is ours: experiments/bnb_straight.py.
%  - Check MOR's policy on AI-assisted research and add the disclosure statement.  [TODO]
\documentclass[11pt]{article}
\usepackage[margin=1in]{geometry}
\usepackage{amsmath,amssymb,amsthm,mathtools,booktabs,enumitem}
\usepackage[hidelinks]{hyperref}
\newtheorem{theorem}{Theorem}
\newtheorem{proposition}[theorem]{Proposition}
\newtheorem{lemma}[theorem]{Lemma}
\newtheorem{corollary}[theorem]{Corollary}
\newtheorem{conjecture}[theorem]{Conjecture}
\theoremstyle{remark}
\newtheorem{remark}[theorem]{Remark}
\newcommand{\R}{\mathbb{R}}
\newcommand{\E}{\mathbb{E}}
\newcommand{\Da}{\mathrm{Da}}
\newcommand{\Pe}{\mathrm{Pe}}
\DeclareMathOperator*{\argmax}{arg\,max}
\DeclareMathOperator*{\argmin}{arg\,min}

\begin{document}
\section{Where should the sensors go? Theory}\label{sec:theory}

The experiments of the previous section found that, among all layouts of six
sensors, the best one is a regular ring whose radius depends on how directed the drone is.
This section states the layout problem precisely, proves what can be proved about it, and
separates those results from the numerical evidence for the rest.

\subsection{The layout problem}
Lengths are measured in sensor ranges ($\rho=1$) and time in range-crossings at full speed
($v=1$). The drone moves in the annulus $\Omega=\{x\in\R^2: d<|x|<R\}$ by
\begin{equation}\label{eq:sde}
dX_t=-\frac{X_t}{|X_t|}\,dt+\sqrt{2D}\,dW_t ,\qquad D=1/\Pe ,
\end{equation}
reflecting at the outer edge $|x|=R$ and arriving when it first reaches the response line
$|x|=d$, at time $\tau$. A layout is $s=(s_1,\dots,s_K)\in(\R^2)^K$, and detection is a
Poisson process with rate
\begin{equation}\label{eq:kappa}
\kappa_s(x)=\Da\sum_{j=1}^K e^{-|x-s_j|^2/2}.
\end{equation}
The chance of arriving undetected from $x$ is
$u_s(x)=\E_x\big[\exp(-\int_0^\tau\kappa_s(X_t)\,dt)\big]$, which by the Feynman--Kac formula
solves
\begin{equation}\label{eq:pde}
D\,\Delta u-\partial_r u-\kappa_s u=0\ \text{in }\Omega,\qquad u=1\ \text{on }|x|=d,\qquad
\partial_r u=0\ \text{on }|x|=R .
\end{equation}
The drone chooses where to enter and the defender chooses the layout:
\begin{equation}\label{eq:minimax}
F(s)=\max_{|x|=R}u_s(x),\qquad F^*=\inf_{s}F(s).
\end{equation}
A \emph{regular ring} of radius $a$ is $s_j=a\,e^{2\pi i j/K}$ (identifying $\R^2$ with
$\mathbb{C}$, up to rotation); we write $F_{\rm ring}(a)$ and $a^*=\argmin_a F_{\rm ring}(a)$.
Throughout, $T_D(x)=\E_x\tau$ is the mean arrival time, which is finite and bounded on
$\overline\Omega$.

\subsection{Basic properties}
\begin{proposition}[Well-posedness]\label{prop:basic}
\begin{enumerate}[label=(\alph*)]
\item \emph{More detection never helps the drone:} if $\kappa_1\le\kappa_2$ then
$u_{\kappa_2}\le u_{\kappa_1}$, strictly on $\overline\Omega\setminus\{|x|=d\}$ if
$\kappa_1<\kappa_2$ on an open set.
\item \emph{Sensors belong inside the area:} replacing a sensor at $s_j=a n$ with $a>R$,
$|n|=1$, by $Rn$ strictly lowers $u_s$ on the outer edge. Hence every minimizer lies in
$(\overline B_R)^K$.
\item \emph{Continuity:} $|F(s)-F(t)|\le \Da\,e^{-1/2}\,\sup_{\overline\Omega}T_D\,
\sum_j|s_j-t_j|$.
\item A minimizer of $F$ exists, and $F$ is invariant under rotations and reflections of the
plane and under relabeling the sensors.
\end{enumerate}
\end{proposition}
\begin{proof}
(a) follows path by path from the Feynman--Kac representation. (b) For $|x|\le R$,
$|x-an|^2-|x-Rn|^2=(a-R)(a+R-2x\cdot n)>0$, so the moved sensor raises $\kappa_s$ everywhere in
$\Omega$; apply (a) and compactness of the outer circle. (c) Coupling both layouts on the
same path, $|e^{-A}-e^{-B}|\le|A-B|\le\tau\|\kappa_s-\kappa_t\|_\infty$, and the gradient of
$e^{-|x-y|^2/2}$ in $y$ has norm at most $e^{-1/2}$. (d) $F$ is continuous on the compact set
$(\overline B_R)^K$, which contains a minimizer by (b).
\end{proof}

\subsection{Rings: where the drone gets through}
For a ring, the detection rate at fixed distance $r$ from the site is
$\kappa(r,\theta)=\Da\,e^{-(r^2+a^2)/2}S_K(ra,\theta)$ with
$S_K(z,\theta)=\sum_{j=0}^{K-1}e^{z\cos(\theta-2\pi j/K)}$.
\begin{lemma}[Coverage is thinnest midway between sensors]\label{lem:midpoint}
For $K\ge2$ and $z>0$, $S_K(z,\theta)$ is strictly decreasing in $\theta$ on $(0,\pi/K)$.
\end{lemma}
\begin{proof}
For $0<\theta<\pi/K$ put $c=\cos K\theta$, which decreases from $1$ to $-1$. The distinct
numbers $t_j=\cos(\theta-2\pi j/K)$ are the roots of $T_K(t)-c$, where $T_K$ is the Chebyshev
polynomial with leading coefficient $2^{K-1}$. Implicit differentiation gives
$dt_j/dc=1/T_K'(t_j)$, so with $f(t)=z e^{zt}$,
\[
\frac{dS_K}{dc}=\sum_j\frac{f(t_j)}{T_K'(t_j)}=2^{1-K}f[t_0,\dots,t_{K-1}]
=\frac{2^{1-K}}{(K-1)!}f^{(K-1)}(\xi)=\frac{2^{1-K}z^K e^{z\xi}}{(K-1)!}>0,
\]
using the divided-difference mean-value theorem. As $c$ decreases with $\theta$, so does $S_K$.
\end{proof}

\begin{theorem}[The drone's best entry point]\label{thm:midpoint}
For every $D>0$ and every ring radius $a>0$, $u$ on the outer edge is strictly increasing in
$\theta$ on $(0,\pi/K)$. Hence $F_{\rm ring}(a)=u(R,\pi/K)$: the drone's best entry point is
always midway between two neighboring sensors.
\end{theorem}
\begin{proof}
Let $w=\partial_\theta u$ on the wedge $W=\{d<r<R,\ 0<\theta<\pi/K\}$. The operator in
\eqref{eq:pde} has coefficients independent of $\theta$, so differentiating gives
$D\Delta w-\partial_r w-\kappa w=(\partial_\theta\kappa)\,u$, which is negative in $W$ by
Lemma~\ref{lem:midpoint} and $u>0$. By the ring's mirror symmetries $w=0$ on $\theta=0$ and
$\theta=\pi/K$; $w=0$ on $r=d$ since $u\equiv1$ there; and $\partial_r w=0$ on $r=R$. The
strong minimum principle (the zeroth-order coefficient $-\kappa$ is nonpositive) and Hopf's
lemma at the Neumann edge exclude a nonpositive minimum of $w$ in $W$ or on its outer arc, so
$w>0$ there.
\end{proof}

\begin{proposition}[The best ring is a stationary point of the full problem]
\label{prop:clarke}
If $a>0$ is a local minimizer of $F_{\rm ring}$, then the ring of radius $a$ is
Clarke-stationary for $F$ on $(\R^2)^K$: $0\in\partial^C F(s(a))$.
\end{proposition}
\begin{proof}
For $D>0$, $u_s(R,\theta)$ is smooth in $s$, so by Danskin's theorem
$C=\partial^CF(s(a))=\operatorname{conv}\{\nabla_s u_s(R,\theta):\theta\in A\}$, where $A$ is
the set of maximizing angles. With $e=s(a)/a$ the radial direction, local minimality along
the ring gives $\max_{g\in C}g\cdot e\ge0$ and $\min_{g\in C}g\cdot e\le0$, so some $h\in C$
has $h\cdot e=0$. The dihedral group $D_K$ acts on layouts by rotating or reflecting the plane
and relabeling sensors; it fixes $s(a)$ and preserves $C$. Its fixed vectors are exactly the
multiples of $e$: rotations force equal displacements in each sensor's own frame, and the
reflection through a sensor removes the tangential part. The group average of $h$ therefore
lies in $C$ and equals $(h\cdot e/K)\,e=0$.
\end{proof}

\subsection{The straight-line drone ($D\to0$)}
As $D\to0$ the drone flies straight in along the ray at its entry angle, and
$u_s(R,\theta)\to e^{-I_s(\theta)}$ with \emph{exposure}
$I_s(\theta)=\int_d^R\kappa_s(re^{i\theta})\,dr$. Minimizing $F$ becomes maximizing
$V(s)=\min_\theta I_s(\theta)$. Because $V$ is proportional to $\Da$, the best layout in this
limit depends only on $(K,d,R)$. One sensor at polar position $(a,\phi)$ contributes the
closed form
\begin{equation}\label{eq:h}
h(a,\theta-\phi)=\sqrt{\tfrac{\pi}{2}}\,e^{-a^2\sin^2 t/2}
\Big[\operatorname{erf}\tfrac{R-a\cos t}{\sqrt2}-\operatorname{erf}\tfrac{d-a\cos t}{\sqrt2}\Big],
\qquad t=\theta-\phi .
\end{equation}
Write $B_n(a)=\int_d^R e^{-(r^2+a^2)/2}I_n(ra)\,dr$ with $I_n$ the modified Bessel function.

\begin{proposition}[Exact ring exposure]\label{prop:ringexposure}
For a ring of radius $a$,
$I_{\rm ring}(\theta)=K\Da\big[B_0(a)+2\sum_{q\ge1}B_{qK}(a)\cos(qK\theta)\big]$, its minimum is
at the midpoints, and
\[
M_K(a):=\min_\theta I_{\rm ring}=K\Da\Big[B_0(a)+2\sum_{q\ge1}(-1)^qB_{qK}(a)\Big],\qquad
0\le K\Da\,B_0(a)-M_K(a)\le 2K\Da\,B_K(a).
\]
\end{proposition}
\begin{proof}
Insert $e^{z\cos\varphi}=I_0(z)+2\sum_{n\ge1}I_n(z)\cos n\varphi$ and sum over the ring; only
multiples of $K$ survive. The minimum location is Lemma~\ref{lem:midpoint} integrated in $r$.
The bound is the alternating-series estimate, since $I_n(z)$ decreases in $n$ for $z>0$.
\end{proof}

\begin{proposition}[What the straight-line problem says]\label{prop:straight}
Let $V^*=\sup_sV(s)$ and $A_0=\int_d^Re^{-r^2/2}dr$, $A_2=\int_d^R(r^2-2)e^{-r^2/2}dr$.
\begin{enumerate}[label=(\alph*)]
\item \emph{Rings are asymptotically optimal:}
$V^*\le K\Da\max_aB_0(a)$, and the ring at $\argmax_a B_0$ is within
$2K\Da(R-d)\,(R^2/2)^K/K!$ of $V^*$.
\item \emph{Collapse to the center:} if $K\le2$, or if $R\le2$ and $A_2\le0$, the unique
optimum puts every sensor at the center.
\item \emph{No collapse:} if $K\ge3$ and $A_2>0$, then
$M_K(a)=K\Da A_0+\tfrac{K\Da}{4}A_2a^2+O(a^3)$, so a small ring beats the center.
\end{enumerate}
\end{proposition}
\begin{proof}[Proof sketch]
(a) $V(s)\le\frac1{2\pi}\int I_s=\Da\sum_jB_0(|s_j|)$; combine with
Proposition~\ref{prop:ringexposure} and $I_K(z)\le (z/2)^Ke^z/K!$.
(b) For $K\le2$ some ray makes a nonacute angle with both sensors, and along it each sensor
contributes at most what a central sensor would. For $R\le2$, $A_2\le0$, the moment bound
$\E_\mu r^{2m}\le2^mm!$ under $\mu\propto e^{-r^2/2}dr$ gives $B_0(a)<A_0$ for $a>0$.
(c) Expand $e^{-a^2/2}I_0(ra)=1+\tfrac{a^2}{4}(r^2-2)+O(a^4)$ and use $B_K=O(a^K)$.
\end{proof}
For the base geometry $d=2$, $R=6.5$ we have $A_2>0$, so collapse never occurs; the bound in
(a) is only useful for large $K$. The straight-line best ring radius is
$\argmax_aM_K(a)$ (for $K=6$: $a=2.4528$, $M_6=1.30410$, $e^{-M_6}=0.2714$), not
$\argmax_aB_0(a)=3.060$, which maximizes \emph{average} rather than worst-case coverage.

\begin{proposition}[Local optimality of the hexagon; computer-assisted]\label{prop:local}
For $K=6$, $d=2$, $R=6.5$, the best ring radius is unique in $[2.452774,\,2.452776]$, and the
regular hexagon there is a strict local maximizer of $V$, modulo rotation.
\end{proposition}
\noindent\emph{Status.} Proved by an interval-arithmetic certificate of a second-order
condition for max-min functions that accounts for movement of the worst entry angle
(see the notes; \textbf{to be re-run from its code before submission}). Independently, all
800 random perturbations of size $3\times10^{-3}$ and $10^{-2}$ that we tried lowered $V$.

\begin{proposition}[Global certificate up to $\varepsilon$; branch and bound]\label{prop:bnb}
For $K=4$, $d=2$, $R=6.5$: no layout has $V(s)>M_4(a^*)+\varepsilon$, with
$M_4(a^*)=0.522161$ at $a^*=1.822662$ and $\varepsilon=2\times10^{-4}$ (0.04\% of the optimum).
\end{proposition}
\noindent\emph{Method.} Branch and bound over the $2K-1$ layout coordinates (after fixing
rotation, ordering the angles and putting the outermost sensor at angle $0$; radii in
$[0,R]$ by Proposition~\ref{prop:basic}(b)). For a box of layouts and finitely many test
angles $\theta$, $\max V\le\min_\theta\sum_j\sup_{\rm box}h$; the supremum over each sensor's
angular range is attained at the closest angle because $h$ decreases in $|t|$, and over its
radial range by tabulated values of \eqref{eq:h} plus the bound
$|\partial_a^2h|\le5.5$. Boxes whose bound is at most the target are discarded; the rest are
halved. Near the ring this first-order bound loses accuracy in proportion to the box width,
so small boxes also use a second-order bound: a Taylor expansion of each sensor's contribution
with a rigorous remainder (Hessian and gradient bounds that shrink with the sensor's distance
from each test ray), combined over test angles by linear-programming dual weights. The final
bound is evaluated from those weights directly, so it is valid whatever the solver's
accuracy. The $K=4$ certificate took $3.2\times10^{6}$ boxes and 56 minutes on one core
(at $\varepsilon=10^{-3}$: $10^5$ boxes, 90 seconds).
For $K=6$ the same method did not finish: after 90 minutes, even with the weak target
$\varepsilon=0.1$, about a quarter of the search space remained open. The bottleneck is the
five free angles. A box that allows each sensor a wide range of angles is bounded as if every
sensor sat at the best angle for each ray, which behaves like a ring smeared into a continuous
band (exposure up to $K\Da\max_aB_0=1.667$, well above the hexagon's $1.304$); only boxes that
are narrow in all five angles can be discarded. Certifying $K=6$ will need a bound that keeps
sensors discrete across rays, or a symmetry reduction, and is left for future work.

\subsection{How the best radius depends on drone randomness}
\begin{proposition}[Both ends of the randomness scale]\label{prop:radius}
\begin{enumerate}[label=(\alph*)]
\item \emph{Nearly straight drones (formal expansion):}
$F_{\rm ring}(a)=e^{-M_K(a)}\big[1+D\,H_K(a)+O(D^2)\big]$ with
\[
H_K(a)=\kappa(d,\tfrac{\pi}{K})+\int_d^R\Big[\kappa^2-\frac{\kappa}{r}
-\frac{J_{\theta\theta}}{r^2}\Big]_{\theta=\pi/K}dr,\qquad
J(r,\theta)=\int_d^r\kappa(q,\theta)\,dq,
\]
the term $\kappa(d,\pi/K)$ coming from the boundary layer at the reflecting edge. Hence
$a^*(D)=a^*(0)+cD+O(D^2)$ with $c=H_K'(a^*(0))/M_K''(a^*(0))$; for the base case
$c\approx3.317$.
\item \emph{Strongly wandering drones:} $u_s=1-w_s/D+O(D^{-2})$ where $-\Delta w_s=\kappa_s$,
$w_s=0$ at $r=d$, $\partial_rw_s=0$ at $r=R$; the best ring radius tends to the maximizer of
$\min_\theta w_{s(a)}(R,\theta)$, which for the base case is $a_\infty\approx4.96$.
\end{enumerate}
\end{proposition}

\begin{table}[h]\centering
\caption{Best ring radius for six sensors ($d=2$, $R=6.5$, $\Da=1$): finite-difference
solutions of \eqref{eq:pde} against the two asymptotic predictions.}\label{tab:radius}
\begin{tabular}{@{}lcccccccccc@{}}\toprule
$D$ & 0 & 0.005 & 0.01 & 0.02 & 0.05 & 0.5 & 2 & 8 & 32 & $\infty$\\\midrule
PDE & 2.453 & 2.471 & 2.487 & 2.517 & 2.599 & 3.25 & 4.03 & 4.67 & 4.89 & \\
$a^*(0)+cD$ & & 2.469 & 2.486 & 2.519 & 2.619 & & & & & \\
Large-$D$ limit & & & & & & & & & & 4.96\\\bottomrule
\end{tabular}
\end{table}
The best radius grows steadily with randomness, from the straight-line value $2.453$ toward
$4.96$: a defender facing wandering drones should push the ring outward, where it covers the
longer, more tortuous approach paths. Monotonicity for all $D$ is supported by
Table~\ref{tab:radius} but not proved.

\subsection{A floor no layout can beat}
\begin{proposition}[Rotational-average lower bound]\label{prop:floor}
For every layout, $F(s)\ge u_{\bar\kappa_s}(R)$, where
$\bar\kappa_s(r)=\Da\sum_je^{-(r^2+|s_j|^2)/2}I_0(r|s_j|)$ is $\kappa_s$ averaged over
rotations. Consequently $F^*\ge\min_{a_1,\dots,a_K\ge0}u_{\bar\kappa(a)}(R)$, a problem with
radial symmetry. If $R\le\sqrt2$ the bound is attained by putting every sensor at the center.
\end{proposition}
\begin{proof}
By Jensen's inequality path by path,
$u_{\bar\kappa_s}(x)\le\frac1{2\pi}\int_0^{2\pi}u_{Q_\phi s}(x)\,d\phi$, and by rotational
invariance of \eqref{eq:sde} the right side at $|x|=R$ is the average of $u_s$ over entry
angles, which is at most $F(s)$. For $R\le\sqrt2$,
$e^{-a^2/2}I_0(ra)\le e^{-a^2/2+r^2a^2/4}\le1$, so $\bar\kappa_s\le K\Da e^{-r^2/2}$.
\end{proof}

\begin{table}[h]\centering
\caption{Six sensors, $d=2$, $R=6.5$, $\Da=1$: the best ring against the floor of
Proposition~\ref{prop:floor} (radial problem solved numerically; the minimizing radii were
all equal).}\label{tab:floor}
\begin{tabular}{@{}lcccc@{}}\toprule
Drone & $\Pe$ & Best ring $F$ & Floor & Gap\\\midrule
Wandering & 0.5 & 0.4440 & 0.4099 & 0.034\\
Mixed & 2 & 0.3101 & 0.2683 & 0.042\\
Directed & 20 & 0.2710 & 0.1986 & 0.072\\\bottomrule
\end{tabular}
\end{table}
The floor is tightest for wandering drones, which effectively average the detection field over
angle. Averaging only over the ring's own $K$-fold rotations does not tighten it: the averaged
field of six suitably placed sensors is that of a 36-sensor ring of one-sixth strength, which
is within $10^{-9}$ of the fully averaged field.

\subsection{Is the ring the best layout? Evidence and conjecture}
\begin{conjecture}\label{conj:ring}
For $d=2$, $R=6.5$, $\Da=1$ and $K=6$, every minimizer of $F$ is a regular ring, for every
$D>0$.
\end{conjecture}
\noindent Evidence, beyond Propositions~\ref{prop:clarke}, \ref{prop:local} and
\ref{prop:bnb}:
\begin{itemize}[nosep]
\item Gradient search (exact adjoint gradients of a soft maximum of the edge values,
L-BFGS-B) from 36 starts per drone type, 30 of them random: all 108 runs ended at a regular
hexagon (spacing $60^\circ$, radii within $0.03$), within $2\times10^{-4}$ of the best ring,
below the grid error.
\item Straight-line limit: 45 epigraph searches (SLSQP), 40 from random layouts, all ended at
the hexagon of radius $2.453$; for $K=4$, likewise at the square of radius $1.823$.
\item Structured families (two staggered rings, zigzag, five plus one, four plus two) either
collapse to the ring or are worse by $0.08$--$0.20$; a seeded evolutionary search never
improved on its ring seed.
\item Pattern changes: for straight-line drones the single ring beats ring-plus-center and
every two-ring split for $K=3,\dots,10$, unlike the small-trap eigenvalue problem, where a
ring plus a center trap wins from seven traps \cite{KTW05}.
\end{itemize}
Conjecture~\ref{conj:ring} cannot hold for every geometry: by
Propositions~\ref{prop:straight}(b) and~\ref{prop:floor}, small areas ($R\le\sqrt2$, or
$R\le2$ with $A_2\le0$ for straight-line drones) are best defended by stacking every sensor at
the center. The ring is the answer when the area is large compared with a sensor's range.

\begin{thebibliography}{9}
\bibitem{KTW05} T.~Kolokolnikov, M.~S.~Titcombe, M.~J.~Ward, Optimizing the fundamental
Neumann eigenvalue for the Laplacian in a domain with small traps, \emph{European J. Appl.
Math.} 16 (2005) 161--200.
\bibitem{Oks} B.~{\O}ksendal, \emph{Stochastic Differential Equations}, 6th ed., Springer,
2003.
\bibitem{Clarke} F.~H.~Clarke, \emph{Optimization and Nonsmooth Analysis}, Wiley, 1983.
\bibitem{GT} D.~Gilbarg, N.~S.~Trudinger, \emph{Elliptic Partial Differential Equations of
Second Order}, Springer, 2001.
\end{thebibliography}
\end{document}
